\section{Implementation}
\subsection{Implementation Variables and Internal Macros}

The package uses a combination of \lstinline|\newcommand| and \lstinline|\def| to implement internal variables and functions. Although both constructs define macros, they serve slightly different purposes and are therefore used for different tasks within the package.

\subsubsection{Internal Variables}

Several package parameters are implemented as macros storing simple values:

\begin{lstlisting}[language={LaTeX}]
\newcommand{\ba@region@scale}{1.25}
\newcommand{\ba@region@radius}{6}
\newcommand{\ba@caption@style}{\scriptsize}
\newcommand{\ba@figure@scale}{0.25}
\newcommand{\ba@caption@shift}{-3mm}
\end{lstlisting}

These macros behave similarly to variables in conventional programming languages. For example,

\begin{lstlisting}[language={LaTeX}]
\ba@region@scale
\end{lstlisting}

expands to

\begin{lstlisting}[language=LaTeX]
1.25
\end{lstlisting}

whenever it is used.

The values may later be modified through the PGF key interface, for example

\begin{lstlisting}[language={LaTeX}]
\baSetup{region scale=1.5}
\end{lstlisting}

which internally redefines the corresponding macro.

\subsubsection{Using \texttt{\textbackslash newcommand}}

The command

\begin{lstlisting}[language={LaTeX}]
\newcommand{\ba@region@scale}{1.25}
\end{lstlisting}

creates a new macro and generates an error if the macro already exists.

This behaviour provides a safety mechanism during package development, since accidental redefinitions are detected immediately.

Consequently, \lstinline|\newcommand| is generally preferred for defining package variables, user commands, and stable interfaces.

\subsubsection{Using \texttt{\textbackslash def}}

The package also contains the definition

\begin{lstlisting}[language={LaTeX}]
\def\ba@fmtcaption{\alph{baSubRegion}}
\end{lstlisting}

which specifies the formatting function used for caption labels.

Unlike \lstinline|\newcommand|, the primitive

\begin{lstlisting}[language={LaTeX}]
\def
\end{lstlisting}

does not check whether a macro already exists and simply replaces its definition.

This makes \lstinline|\def| useful when a macro is intended to be repeatedly overwritten by internal package logic. In the present package, the caption format may change dynamically according to the selected key

\begin{lstlisting}[language={LaTeX}]
caption label=alpha
caption label=roman
caption label=Roman
caption label=arabic
\end{lstlisting}

which internally execute statements such as

\begin{lstlisting}[language={LaTeX}]
\def\ba@fmtcaption{\roman{baSubRegion}}
\end{lstlisting}

or

\begin{lstlisting}[language={LaTeX}]
\def\ba@fmtcaption{\arabic{baSubRegion}}
\end{lstlisting}

to replace the current formatting rule.

For this reason, \lstinline|\def| is often used for internal implementation details whose meaning changes dynamically.

\subsubsection{Variable versus Function Macros}

A useful distinction is to regard

\begin{lstlisting}[language={LaTeX}]
\newcommand{\ba@region@scale}{1.25}
\end{lstlisting}

as a variable storing a value, whereas

\begin{lstlisting}[language={LaTeX}]
\def\ba@fmtcaption{\alph{baSubRegion}}
\end{lstlisting}

acts more like a function that computes a representation from the current counter value.

The user-level command

\begin{lstlisting}[language={LaTeX}]
\newcommand{\baFormatCaptionLabel}{\ba@fmtcaption}
\end{lstlisting}

provides an abstraction layer. All package code accesses caption labels through

\begin{lstlisting}[language={LaTeX}]
\baFormatCaptionLabel
\end{lstlisting}

without needing to know which numbering scheme is currently active.

\subsubsection{The Meaning of the \texttt{@} Character}

Many internal macros contain the character \lstinline|@|:

\begin{lstlisting}[language={LaTeX}]
\ba@region@scale
\ba@region@radius
\ba@fmtcaption
\ba@atlaslabel
\end{lstlisting}

In LaTeX, \lstinline|@| is traditionally used to distinguish internal implementation macros from user-level commands.

When a package is loaded, \lstinline|@| is temporarily treated as a letter, allowing definitions such as

\begin{lstlisting}[language={LaTeX}]
\newcommand{\ba@region@scale}{1.25}
\end{lstlisting}

to be created.

Outside package code, however, \lstinline|@| is not considered a letter and therefore these commands are normally inaccessible to users.

This convention serves two purposes:

\begin{enumerate}
\item It separates the public interface from the internal implementation.
\item It reduces the risk of name collisions with user-defined commands.
\end{enumerate}

For example,

\begin{lstlisting}[language={LaTeX}]
\baSetup
\baDeclareRegion
\baAtlasFigure
\end{lstlisting}

form part of the public API, whereas

\begin{lstlisting}[language={LaTeX}]
\ba@region@scale
\ba@fmtcaption
\ba@atlaslabel
\end{lstlisting}

are intended solely for internal package use.

The naming convention therefore clearly distinguishes package internals from user-accessible commands and improves maintainability of the implementation.

\subsection{Macros and Variables in \LaTeX}

In \LaTeX{}, the concepts of \emph{variables} and \emph{macros} do not exist as separate language constructs in the same way as in programming languages such as Python or C. Instead, both concepts are implemented using \emph{macros}. Understanding this distinction is essential for designing robust package code.

\subsubsection{Macros as the Fundamental Mechanism}

A macro in \LaTeX{} is a named command that expands into replacement text. For example:

\begin{lstlisting}[language=LaTeX]
\newcommand{\example}{Hello World}
\end{lstlisting}
Whenever \lstinline|\example| is used, it is replaced by its definition:

\begin{lstlisting}[language=LaTeX]
Hello World
\end{lstlisting}
Macros can therefore represent:

\begin{itemize}
\item text fragments,
\item formatting instructions,
\item mathematical expressions,
\item or even code-like behavior.
\end{itemize}
In this sense, macros are the \emph{fundamental abstraction mechanism} of \TeX{}.

\subsubsection{Variables as a Conceptual Interpretation}

Although \LaTeX{} does not provide true variables, macros are often used to simulate them. For example:
\begin{lstlisting}[language=LaTeX]
\newcommand{\ba@region@scale}{1.25}
\end{lstlisting}
can be interpreted as a \emph{variable} storing a numeric value.

However, technically it remains a macro that expands to the token sequence \lstinline|1.25|. There is no memory cell, no type system, and no assignment in the traditional sense.

Thus, a ``variable'' in \LaTeX{} is simply a macro whose expansion is treated as data.

\subsubsection{Difference Between Macros and Variables (Conceptual View)}

The distinction is therefore conceptual rather than technical:

\begin{center}
\begin{tabular}{ll}
\textbf{Macro} & \textbf{Variable (conceptual)} \\
\hline
General-purpose expansion rule & Stores a value-like token sequence \\
May contain logic or formatting & Intended to behave like data \\
Can act as function or constant & Typically used as a constant \\
\end{tabular}
\end{center}
For example:

\begin{lstlisting}[language=LaTeX]
\def\ba@fmtcaption{\alph{baSubRegion}}
\end{lstlisting}
behaves like a \emph{function macro}, since it computes an output based on a counter.

In contrast:

\begin{lstlisting}[language=LaTeX]
\newcommand{\ba@region@radius}{6}
\end{lstlisting}
behaves like a \emph{variable macro}, since it stores a fixed value.

\subsubsection{Assignment and Mutability}

Unlike variables in imperative programming languages, macros do not support assignment. Instead, they are redefined:

\begin{lstlisting}[language=LaTeX]
\renewcommand{\ba@region@radius}{7}
\end{lstlisting}
or

\begin{lstlisting}[language=LaTeX]
\def\ba@fmtcaption{\roman{baSubRegion}}
\end{lstlisting}

This means that what appears to be ``updating a variable'' is actually replacing one macro definition with another.

\subsubsection{Summary}

In summary, \LaTeX{} does not distinguish between macros and variables at the language level. Instead:

\begin{itemize}
\item \textbf{Macros} are the underlying mechanism for all definitions.
\item \textbf{Variables} are a conceptual interpretation of macros used as data containers.
\item \textbf{Functions} are macros whose expansion depends on arguments or counters.
\end{itemize}

This abstraction allows flexible package design but requires careful naming and structure to avoid unintended redefinitions or expansion issues.

\subsection{Floating Environment and Atlas State Management}

The \texttt{bifurcationatlas} package builds on the \texttt{newfloat} mechanism to provide a structured floating environment for bifurcation diagrams, while maintaining an internal state system for consistent numbering, labeling, and region management.

\subsubsection{Floating Environment Declaration}

The core floating container is defined using

\begin{lstlisting}[language=LaTeX]
\DeclareFloatingEnvironment[
  fileext=loa,
  listname={List of Bifurcation Atlases},
  name=Bifurcation Atlas,
  placement=htbp
]{bifurcationatlasfloat}
\end{lstlisting}

This creates a new float type similar to \texttt{figure} or \texttt{table}. The configuration has the following effects:

\begin{itemize}
\item \texttt{fileext=loa}: writes entries to a dedicated list file (analogous to \texttt{lof} or \texttt{lot}).
\item \texttt{listname}: defines the title of the generated list of atlases.
\item \texttt{name}: sets the printed caption prefix.
\item \texttt{placement=htbp}: controls float placement rules.
\end{itemize}

This float acts as the visual and structural container for each bifurcation atlas instance.

\subsubsection{Global Atlas State}

The package maintains a minimal global state to ensure consistency across nested macros:

\begin{lstlisting}[language=LaTeX]
\newif\ifba@inatlas
\def\ba@atlaslabel{}
\let\ba@oldlabel\label
\end{lstlisting}
\paragraph{The Syntax and Meaning of \texttt{\textbackslash newif}}

The primitive \lstinline|\newif| is a low-level \TeX{} construct used to define conditional flags (booleans). Its syntax

\begin{lstlisting}[language=LaTeX]
\newif\ifba@inatlas
\end{lstlisting}

creates a new conditional named \lstinline|\ifba@inatlas| together with two associated commands:

\begin{itemize}
\item \lstinline|\ba@inatlastrue| sets the conditional to \emph{true},
\item \lstinline|\ba@inatlasfalse| sets the conditional to \emph{false}.
\end{itemize}

After declaration, the conditional can be used in standard \TeX{} branching structures:

\begin{lstlisting}[language=LaTeX]
\ifba@inatlas
  % code executed if true
\else
  % code executed if false
\fi
\end{lstlisting}

Unlike higher-level constructs such as \lstinline|\newcommand| or \LaTeX3 booleans, \lstinline|\newif| operates at the primitive \TeX{} level and provides a lightweight mechanism for global state management.

A key property is that the conditional itself (\lstinline|\ifba@inatlas|) is not a variable in the usual sense, but a control-flow switch that expands directly into branching logic during compilation.

In package design, such conditionals are typically used to track execution context, for example whether code is currently inside a specific environment, ensuring that certain commands are only valid in the correct scope.
\begin{itemize}
\item \lstinline|\ifba@inatlas| is a boolean flag indicating whether the current code is executed inside an atlas environment.
\item \lstinline|\ba@atlaslabel| stores the user-defined label of the current atlas (if provided via \lstinline|\label|).
\item \lstinline|\ba@oldlabel| preserves the original LaTeX labeling mechanism before it is temporarily overridden.
\end{itemize}
Additionally, the macro

\begin{lstlisting}[language=LaTeX]
\newcommand{\baRequireAtlas}{...}
\end{lstlisting}

acts as a safety guard. It produces a package error if atlas-specific commands are used outside the environment, preventing inconsistent state usage.

\subsubsection{Atlas Environment Initialization}

The environment is defined using

\begin{lstlisting}[language=LaTeX]
\NewDocumentEnvironment{bifurcationatlas}{O{}}{ ... }{ ... }
\end{lstlisting}

where the optional argument allows configuration via a key--value interface.

At the beginning of the environment, a local group is opened:

\begin{lstlisting}[language=LaTeX]
\begingroup
\end{lstlisting}

This ensures that all internal changes remain local and do not leak into the global document state.

\paragraph{User Configuration}

User-provided options are processed via:

\begin{lstlisting}[language=LaTeX]
\pgfkeys{/ba,#1}
\end{lstlisting}
\paragraph{Syntax and Meaning of \texttt{\textbackslash pgfkeys\{/ba,\#1\}}}

The command \lstinline|\pgfkeys{/ba,#1}| is part of the \texttt{pgfkeys} key--value processing system and is used to apply a set of user-provided configuration options to a predefined key family.

The syntax consists of two components:

\begin{itemize}
\item \texttt{/ba} is the key family (or key path prefix), defining the namespace in which all relevant keys are registered (e.g., \texttt{/ba/region scale}, \texttt{/ba/caption label}).
\item \lstinline|#1| is the argument passed to the environment or macro, containing a comma-separated list of user options such as \texttt{region scale=1.5, caption label=roman}.
\end{itemize}

When executed, \lstinline|\pgfkeys| parses the input string and dispatches each key--value pair to its corresponding handler defined in the key family. For example, an input

\begin{lstlisting}[language=LaTeX]
\begin{bifurcationatlas}[region scale=1.5, caption label=roman]
\end{lstlisting}

is internally expanded into:

\begin{lstlisting}[language=LaTeX]
\pgfkeys{/ba, region scale=1.5, caption label=roman}
\end{lstlisting}

Each key is then evaluated according to its definition in the \texttt{/ba} key family, typically using handlers such as \lstinline|.code|, \lstinline|.store in|, or \lstinline|.is choice|.

Thus, \lstinline|\pgfkeys{/ba,#1}| acts as the central interface between user-level configuration and internal package parameters, enabling structured and extensible option handling without requiring manual parsing of optional arguments.

This connects the environment directly to the package key system, allowing local overrides of visual and structural parameters.

\paragraph{Activation and Counter Reset}

The atlas state is activated and counters are initialized:

\begin{lstlisting}[language=LaTeX]
\ba@inatlastrue
\stepcounter{baAtlasID}
\setcounter{baRegion}{0}
\setcounter{baSubRegion}{0}
\end{lstlisting}

This ensures that each atlas instance receives a unique identifier, while its internal region numbering starts from zero.

\paragraph{Local Storage Reset}

Internal accumulators used during rendering are cleared:

\begin{lstlisting}[language=LaTeX]
\def\ba@regionlist{}
\def\ba@colors{}
\def\ba@caption{}
\end{lstlisting}

These macros act as temporary storage containers for region metadata and are reset for each atlas.

\subsubsection{Safe Label Capture Mechanism}

A central feature of the implementation is the controlled interception of \lstinline|\label|:

\begin{lstlisting}[language=LaTeX]
\let\ba@oldlabel\label
\def\ba@atlaslabel{}
\renewcommand{\label}[1]{%
  \ifx\ba@atlaslabel\@empty
    \gdef\ba@atlaslabel{##1}%
  \fi
  \ba@oldlabel{##1}%
}
\end{lstlisting}
\paragraph{Why \texttt{\#\#1} is required instead of \texttt{\#1}}

Inside the redefinition of a macro, TeX distinguishes between parameter placeholders of the \emph{outer macro being defined} and literal parameter tokens that should be stored inside the newly defined macro. In the construct

\begin{lstlisting}[language=LaTeX]
\renewcommand{\label}[1]{%
  \gdef\ba@atlaslabel{##1}%
}
\end{lstlisting}

the command \lstinline|\label| itself has one argument, denoted by \lstinline|#1|. However, the body of the definition is processed at a higher syntactic level: it is not yet executed but instead being written as the replacement text of \lstinline|\label|. For this reason, any parameter marker that should survive into the final definition must be escaped by doubling the hash symbol.

Thus, \lstinline|##1| does not refer to a second argument, but rather encodes a literal \lstinline|#1| that will only be interpreted when \lstinline|\label| is later executed. If one were to write \lstinline|#1| instead, TeX would attempt to interpret it as a parameter of the surrounding definition process, leading to an ``Illegal parameter number'' error.

In summary, \lstinline|##1| is required to preserve the intended argument reference inside a macro definition, ensuring correct expansion at execution time rather than at definition time.

This mechanism serves two purposes:

\begin{itemize}
\item It captures the first user-provided label as the atlas identifier.
\item It preserves the original \LaTeX{} labeling functionality via \lstinline|\ba@oldlabel|.
\end{itemize}

The redefined \lstinline|\label| macro stores the first encountered label in \lstinline|\ba@atlaslabel| and then delegates to the original implementation. This ensures compatibility with \texttt{hyperref} and standard referencing mechanisms.

\subsubsection{Detailed Execution Logic}

When \lstinline|\label{key}| is executed inside the environment, the macro expansion proceeds as follows:

\begin{enumerate}
\item The argument \lstinline|key| is passed to the redefined \lstinline|\label|.
\item A check is performed:
\begin{lstlisting}[language=LaTeX]
\ifx\ba@atlaslabel\@empty
\end{lstlisting}
This verifies whether an atlas identifier has already been stored.
\item If no identifier exists, the first encountered label is globally stored:
\begin{lstlisting}[language=LaTeX]
\gdef\ba@atlaslabel{key}
\end{lstlisting}
The use of \lstinline|\gdef| ensures persistence across local grouping boundaries and nested environments.
\item The original labeling mechanism is executed:
\begin{lstlisting}[language=LaTeX]
\ba@oldlabel{key}
\end{lstlisting}
This guarantees full compatibility with \LaTeX{} cross-referencing and hyperlinking systems.
\end{enumerate}
\subsubsection{Why Redefining \texttt{\textbackslash label} is Necessary}

The redefinition of \lstinline|\label| is not required to preserve its original functionality, but rather to extend it with additional structural information needed by the bifurcation atlas system. In standard \LaTeX{}, the command \lstinline|\label| only writes an entry to the auxiliary file and enables later referencing via \lstinline|\ref| or \lstinline|\pageref|. However, it does not expose the label key (e.g.\ \texttt{ba:test1}) in a way that can be programmatically reused during compilation.

By redefining \lstinline|\label|, the system introduces a controlled interception mechanism:

\begin{itemize}
\item It detects when a label is issued inside the atlas environment.
\item It captures the first encountered label as the atlas identifier via \lstinline|\ba@atlaslabel|.
\item It preserves all standard \LaTeX{} functionality by delegating execution to the original implementation stored in \lstinline|\ba@oldlabel|.
\end{itemize}

The key idea is that the original behavior is not replaced but wrapped. The line

\begin{lstlisting}[language=LaTeX]
\ba@oldlabel{##1}
\end{lstlisting}

ensures full compatibility with \texttt{hyperref}, auxiliary file writing, and cross-referencing mechanisms.

This design is necessary because \LaTeX{} does not provide a native mechanism to query or structure labels after they are defined. The atlas system therefore uses \lstinline|\label| as a data source, extracting semantic information (the atlas identifier) while maintaining full backward compatibility with the standard referencing system.
\subsubsection{Design Interpretation}

This construction effectively turns \lstinline|\label| into a \emph{dual-purpose hook}:

\begin{itemize}
\item a metadata extractor (capturing the atlas identifier),
\item a transparent pass-through to the original referencing system.
\end{itemize}

Importantly, only the first invocation of \lstinline|\label| has semantic significance for the atlas structure. Subsequent labels behave exactly as in standard \LaTeX{}, ensuring no interference with user-defined cross-references.

\subsubsection{Robustness Considerations}

The mechanism relies on the invariant that \lstinline|\ba@atlaslabel| is initially empty. This avoids ambiguity between an undefined macro and an intentionally empty state. Furthermore, the use of \lstinline|\let| to store \lstinline|\ba@oldlabel| is essential to prevent infinite recursion, since \lstinline|\renewcommand{\label}| would otherwise overwrite the original definition.
\subsubsection{Floating Container Opening}

The actual visual content is placed inside the float:

\begin{lstlisting}[language=LaTeX]
\begin{bifurcationatlasfloat}
\end{lstlisting}

At this stage, all atlas content (background figures, radial layouts, and region nodes) is inserted.

\subsubsection{Environment Closure and Cleanup}

At the end of the environment, the float is closed:

\begin{lstlisting}[language=LaTeX]
\end{bifurcationatlasfloat}
\end{lstlisting}

Afterwards, the system performs cleanup:

\begin{lstlisting}[language=LaTeX]
\ba@inatlasfalse
\let\label\ba@oldlabel
\endgroup
\end{lstlisting}

This restores:

\begin{itemize}
\item the atlas state flag,
\item the original \LaTeX\ label definition,
\item and the previous grouping context.
\end{itemize}

\subsubsection{Summary of Design Principles}

This implementation follows three key design principles:

\begin{itemize}
\item \textbf{Isolation}: Each atlas is fully self-contained via grouping.
\item \textbf{Compatibility}: Standard LaTeX labeling and hyperref mechanisms are preserved.
\item \textbf{Safety}: Misuse outside the environment is explicitly detected.
\end{itemize}

The result is a floating environment that behaves like a standard LaTeX float while supporting structured internal indexing and safe label management.

\subsection{Differences Between \texttt{\textbackslash def}, \texttt{\textbackslash let}, \texttt{\textbackslash gdef}, and \texttt{\textbackslash edef}}

In \TeX{} programming, macro assignment is controlled by a small set of primitive operations that differ significantly in scope, expansion behavior, and evaluation timing. The most important distinctions arise between \lstinline|\def|, \lstinline|\let|, \lstinline|\gdef|, and \lstinline|\edef|.

\subsubsection{Local Definition: \texttt{\textbackslash def}}

The primitive

\begin{lstlisting}[language=LaTeX]
\def\macro{replacement text}
\end{lstlisting}

defines a macro by assigning a replacement text that is expanded when the macro is used. If the macro already exists, it is overwritten.

A key property is that \lstinline|\def| respects grouping: changes made inside \lstinline|\begingroup...\endgroup| remain local and are discarded afterwards.

Thus, \lstinline|\def| is typically used for configurable or temporary definitions within a local scope.

\subsubsection{Direct Assignment: \texttt{\textbackslash let}}

The primitive

\begin{lstlisting}[language=LaTeX]
\let\macroA=\macroB
\end{lstlisting}

assigns \lstinline|\macroA| to be an exact copy of the current meaning of \lstinline|\macroB|. Unlike \lstinline|\def|, it does not create a new replacement text but instead binds the control sequence directly.

Important properties:

\begin{itemize}
\item It is instantaneous (no expansion of content).
\item It creates a fixed alias (snapshot semantics).
\item It is commonly used to preserve original definitions (e.g., saving \lstinline|\label| before redefining it).
\end{itemize}

\subsubsection{Global Definition: \texttt{\textbackslash gdef}}

The primitive

\begin{lstlisting}[language=LaTeX]
\gdef\macro{replacement text}
\end{lstlisting}

is equivalent to \lstinline|\def| but ignores grouping boundaries. This means:

\begin{itemize}
\item The definition is global and persists beyond any \lstinline|\begingroup|.
\item It overrides any local scope restrictions.
\item It is typically used when state must survive environment boundaries.
\end{itemize}

Because of its non-local behavior, \lstinline|\gdef| must be used carefully in package design, especially when environments rely on grouping for cleanup.

\subsubsection{Expansion Definition: \texttt{\textbackslash edef}}

The primitive

\begin{lstlisting}[language=LaTeX]
\edef\macro{replacement text}
\end{lstlisting}

performs a fully expanded definition: the right-hand side is expanded immediately at definition time, and the resulting fully expanded text is stored as the macro body.

Key properties:

\begin{itemize}
\item All expandable tokens are expanded immediately.
\item The final macro contains only unexpandable primitives.
\item It is useful for constructing fixed strings, labels, or keys.
\end{itemize}

However, care must be taken because unintended expansion of fragile commands can lead to errors.

\subsubsection{Fully Expanding Assignment with \texttt{\textbackslash xdef}}

The command

\begin{lstlisting}[language=LaTeX]
\xdef\ba@c\i{\c}
\end{lstlisting}

appears inside constructions such as:

\begin{lstlisting}[language=LaTeX]
\foreach [count=\i] \c in \colors {%
  \expandafter\xdef\csname ba@c\i\endcsname{\c}%
}
\end{lstlisting}

and is a key mechanism for creating indexed macros.



\paragraph{Meaning of \texttt{\textbackslash xdef}}

The command \verb|\xdef| is short for:

\begin{center}
\textbf{e\textbf{x}pandable \textbf{def}inition}
\end{center}

It behaves like:

\begin{itemize}
\item \verb|\def| (defines a macro),
\item but with \textbf{full expansion of the right-hand side},
\item and the definition is made \textbf{global}.
\end{itemize}



\paragraph{Difference to \texttt{\textbackslash def}}

\begin{itemize}
\item \verb|\def| stores the argument \emph{as-is} (no full expansion).
\item \verb|\xdef| fully expands all expandable content before storing it.
\end{itemize}

For example:

\begin{lstlisting}[language=LaTeX]
\def\a{\textbf{red}}
\xdef\b{\a}
\end{lstlisting}

results in:

\begin{itemize}
\item \verb|\a| expands later to \textbf{red}
\item \verb|\b| already contains \textbf{red} at definition time
\end{itemize}



\paragraph{Why \texttt{\textbackslash xdef} is used in indexed color storage}

In the expression:

\begin{lstlisting}[language=LaTeX]
\xdef\csname ba@c\i\endcsname{\c}
\end{lstlisting}

the purpose is to:

\begin{enumerate}
\item Create a macro name dynamically (e.g., \verb|\ba@c1|, \verb|\ba@c2|, ...),
\item Store the actual color value inside it,
\item Ensure the value is fully resolved at assignment time.
\end{enumerate}

This guarantees that no further expansion is required during rendering.



\paragraph{Role of \texttt{\textbackslash expandafter}}

Because \verb|\csname...\endcsname| must be constructed before \verb|\xdef| sees it, we use:

\begin{lstlisting}[language=LaTeX]
\expandafter\xdef\csname ba@c\i\endcsname{\c}
\end{lstlisting}

This forces TeX to first evaluate:

\begin{lstlisting}[language=LaTeX]
\csname ba@c\i\endcsname
\end{lstlisting}

before applying \verb|\xdef|.



\paragraph{Why global assignment matters}

Since region rendering and color lookup happen later in possibly different grouping levels (e.g., inside TikZ pictures or floats), the assignment must persist globally. \verb|\xdef| ensures:

\begin{itemize}
\item values are not lost at group boundaries,
\item indexed macros remain accessible throughout the document,
\item consistency across multiple atlas environments.
\end{itemize}



\paragraph{Summary}

\begin{itemize}
\item \verb|\def|: local, non-expanding assignment
\item \verb|\xdef|: global, fully expanding assignment
\item used here for building a persistent indexed color table
\end{itemize}

\subsubsection{Comparison of the Four Primitives}

\begin{center}
\begin{tabular}{l|l|l|l}
\textbf{Command} & \textbf{Scope} & \textbf{Expansion} & \textbf{Purpose} \\
\hline
\lstinline|\def|  & Local (group-aware) & No immediate expansion & Standard macro definition \\
\lstinline|\let|  & Local/global (copy)  & No expansion           & Alias / preservation \\
\lstinline|\gdef| & Global               & No immediate expansion & Persistent state \\
\lstinline|\edef| & Local/global         & Full expansion         & Precomputed fixed content \\
\lstinline|\xdef| & Global               & Full expansion         & Global precomputed definition \\
\end{tabular}
\end{center}

\subsubsection{Implications for Package Design}

In structured packages such as \texttt{bifurcationatlas}, each primitive serves a distinct role:

\begin{itemize}
\item \lstinline|\def| is used for configurable parameters (e.g., formatting rules).
\item \lstinline|\let| is used to preserve original commands (e.g., saving \lstinline|\label|).
\item \lstinline|\gdef| is used when state must persist across environments (e.g., capturing the atlas label globally).
\item \lstinline|\edef| is used when constructing deterministic keys or auxiliary file entries.
\end{itemize}

Understanding these differences is essential for avoiding subtle bugs related to scope leakage, premature expansion, or unintended overwriting of core \LaTeX{} mechanisms.
\subsection{The Meaning and Syntax of \texttt{\textbackslash label}}

In \LaTeX{}, the command \lstinline|\label| is the fundamental mechanism for creating cross-references. It is used to associate a symbolic name with the current value of a counter or a structural position in the document (such as a section, figure, equation, or custom environment).

\subsubsection{Original Meaning of \texttt{\textbackslash label}}

Internally, \lstinline|\label| does not directly store the visible number shown in the document. Instead, it writes an entry into the auxiliary file (\texttt{.aux}), linking a user-defined key to several pieces of internal information, such as:

\begin{itemize}
\item the current value of the most recently stepped counter,
\item the page number,
\item and additional metadata used by packages such as \texttt{hyperref}.
\end{itemize}

For example, when LaTeX processes:

\begin{lstlisting}[language=LaTeX]
\section{Introduction}\label{sec:intro}
\end{lstlisting}

it writes an entry similar to:

\begin{lstlisting}[language=LaTeX]
\newlabel{sec:intro}{{1}{1}}
\end{lstlisting}

This means that the label \lstinline|sec:intro| refers to section number 1 on page 1.

\subsubsection{Syntax of \texttt{\textbackslash label}}

The general syntax is:

\begin{lstlisting}[language=LaTeX]
\label{key}
\end{lstlisting}

where:

\begin{itemize}
\item \texttt{key} is a user-defined identifier,
\item it is later used with \lstinline|\ref{key}| or \lstinline|\pageref{key}|.
\end{itemize}

The key must be unique within the document to avoid ambiguous references.

\subsubsection{How \texttt{\textbackslash label} Works Internally}

The command \lstinline|\label| stores information by writing to the auxiliary file:

\begin{lstlisting}[language=LaTeX]
\newlabel{key}{{current-number}{page}}
\end{lstlisting}

When the document is compiled again, LaTeX reads this file and resolves all references.

With the \texttt{hyperref} package, the structure is extended to include hyperlink targets:

\begin{lstlisting}[language=LaTeX]
\newlabel{key}{{current-number}{page}{text}{anchor}{}}
\end{lstlisting}

This allows clickable references in the compiled PDF.

\subsubsection{Relation to Counters}

A crucial property of \lstinline|\label| is that it does not generate numbers itself. Instead, it captures the value of the most recently incremented counter.

For example:

\begin{lstlisting}[language=LaTeX]
\refstepcounter{section}
\label{sec:example}
\end{lstlisting}

ensures that the label is associated with the current section counter value.

This is why environments such as \texttt{figure}, \texttt{equation}, and custom structures must use \lstinline|\refstepcounter| to make labels meaningful.

\subsubsection{Usage in Custom Environments}

In custom packages, \lstinline|\label| is often redefined or intercepted to attach additional structure to references. However, its original role remains:

\begin{quote}
to bind a symbolic name to the current counter state and page position for later retrieval via \texttt{\textbackslash ref}.
\end{quote}

To preserve compatibility with standard LaTeX and packages such as \texttt{hyperref}, any redefinition should always delegate to the original \lstinline|\label| using a stored copy (e.g., via \lstinline|\let\oldlabel\label|).

\subsubsection{Summary}

\begin{itemize}
\item \texttt{\textbackslash label} creates a symbolic reference key.
\item It stores counter and page information in the auxiliary file.
\item It does not generate numbers itself but captures existing counter states.
\item It is resolved on the next compilation run via \texttt{\textbackslash ref}.
\item It is the foundation of all cross-referencing mechanisms in \LaTeX.
\end{itemize}
\subsection{Key--Value Interface with \texttt{pgfkeys}: Syntax, Semantics, and Application}

The \texttt{pgfkeys} system (from the PGF/TikZ ecosystem) provides a structured key--value interface for configuring packages and environments in a scalable and hierarchical way. It replaces ad-hoc optional arguments with a declarative configuration model, where user options are parsed as named keys rather than positional parameters.

\subsubsection{General Syntax of \texttt{pgfkeys}}

The basic syntax is:

\begin{lstlisting}[language=LaTeX]
\pgfkeys{<key path>/<key>/.code = {<code>}}
\end{lstlisting}

or equivalently:

\begin{lstlisting}[language=LaTeX]
\pgfkeys{<family>, <key>=<value>, ...}
\end{lstlisting}

Each key is associated with an action, such as:

\begin{itemize}
\item executing code (\texttt{.code}),
\item storing values (\texttt{.store in}),
\item defining choices (\texttt{.is choice}),
\item or nesting into subtrees (hierarchical keys).
\end{itemize}



\subsubsection{Key Family Definition}

In the presented implementation, all keys are grouped under the family:

\begin{lstlisting}[language=LaTeX]
/ba/.is family,
\end{lstlisting}

This establishes a namespace \texttt{/ba} that prevents collisions with other packages. Any key defined under this family must be prefixed by \texttt{/ba/}.

The line

\begin{lstlisting}[language=LaTeX]
/ba,
\end{lstlisting}

sets \texttt{/ba} as the default path, so that keys can be written without repeating the full prefix.



\subsubsection{Direct Code Keys (\texttt{.code})}

Several keys directly execute \LaTeX{} code when used:

\begin{lstlisting}[language=LaTeX]
region scale/.code   = {\renewcommand{\ba@region@scale}{#1}},
region radius/.code  = {\renewcommand{\ba@region@radius}{#1}},
caption style/.code  = {\renewcommand{\ba@caption@style}{#1}},
figure scale/.code   = {\renewcommand{\ba@figure@scale}{#1}},
caption shift/.code  = {\renewcommand{\ba@caption@shift}{#1}},
\end{lstlisting}
These keys map user input directly to internal macros. For example:

\begin{lstlisting}[language=LaTeX]
\baSetup{region radius=7}
\end{lstlisting}
will execute:
\begin{lstlisting}[language=LaTeX]
\renewcommand{\ba@region@radius}{7}
\end{lstlisting}
This approach allows dynamic reconfiguration of rendering parameters without exposing internal macros to the user.

\subsubsection{Value Storage Keys (\texttt{.store in})}

Some keys store values in internal macros:
\begin{lstlisting}[language=LaTeX]
internal/region colors/.store in=\ba@colors,
internal/region caption/.store in=\ba@caption,
\end{lstlisting}
This mechanism assigns the user-provided value directly to a macro. For example:
\begin{lstlisting}[language=LaTeX]
internal/region colors={red,blue,green}
\end{lstlisting}
results in:
\begin{lstlisting}[language=LaTeX]
\def\ba@colors{red,blue,green}
\end{lstlisting}
This is particularly useful for data-like configuration such as lists or structured input.

\subsubsection{Choice Keys (\texttt{.is choice})}

The caption formatting system is implemented as a choice key:

\begin{lstlisting}[language=LaTeX]
caption label/.is choice,
\end{lstlisting}
with variants:
\begin{lstlisting}[language=LaTeX]
/ba/caption label/alpha/.code  = {\def\ba@fmtcaption{\alph{baSubRegion}}},
/ba/caption label/roman/.code  = {\def\ba@fmtcaption{\roman{baSubRegion}}},
/ba/caption label/Roman/.code  = {\def\ba@fmtcaption{\Roman{baSubRegion}}},
/ba/caption label/arabic/.code = {\def\ba@fmtcaption{\arabic{baSubRegion}}},
\end{lstlisting}
This defines a controlled enumeration of valid formatting modes. The user can select a mode via:
\begin{lstlisting}[language=LaTeX]
\baSetup{caption label=roman}
\end{lstlisting}
Internally, this sets the macro \lstinline|\ba@fmtcaption| accordingly, which is later used to format region labels.

\subsubsection{User Interface Command: \texttt{\textbackslash baSetup}}

To provide a clean interface, the package defines:

\begin{lstlisting}[language=LaTeX]
\newcommand{\baSetup}[1]{\pgfkeys{/ba,#1}}
\end{lstlisting}

This command forwards all user input into the \texttt{/ba} key family. Thus:

\begin{lstlisting}[language=LaTeX]
\baSetup{region radius=6, caption label=Roman}
\end{lstlisting}

is equivalent to manually calling:

\begin{lstlisting}[language=LaTeX]
\pgfkeys{/ba/region radius=6, /ba/caption label=Roman}
\end{lstlisting}



\subsubsection{Role in the Package Architecture}

Within the \texttt{bifurcationatlas} system, \texttt{pgfkeys} serves as the primary configuration layer. It cleanly separates:

\begin{itemize}
\item \textbf{User interface}: declarative key--value options,
\item \textbf{Internal state}: macros such as \lstinline|\ba@region@radius|,
\item \textbf{Rendering logic}: TikZ-based layout and drawing.
\end{itemize}

This design ensures extensibility, since new options can be added simply by defining additional keys without modifying the core environment logic.

\subsection{Region Registry, Dynamic Macros, and Color Resolution Mechanism}

This section explains the internal design of the region system, which combines dynamic macro creation, list processing, and TikZ style registration. The implementation is based on three core mechanisms: (i) region declaration, (ii) safe access via conditional macros, and (iii) indexed color resolution.

% =========================================================
\subsubsection{Region Declaration System}
% =========================================================

The region declaration command is defined as:

\begin{lstlisting}[language=LaTeX]
\NewDocumentCommand{\baDeclareRegion}{O{} m m}{%
  \baRequireAtlas
  \gappto\ba@regionlist{#2,}
  \expandafter\gdef\csname ba@caption@#2\endcsname{#1}

  \tikzset{
    ba/region/#2/.style={
      /ba/internal/region colors={#3}
    }
  }
}
\end{lstlisting}

\paragraph{Syntax breakdown}

The command has the signature:

\begin{itemize}
\item \texttt{O\{\}}: optional caption text
\item \texttt{m}: region identifier
\item \texttt{m}: region-specific color list
\end{itemize}

\paragraph{Step 1: Environment safety check}

\begin{lstlisting}[language=LaTeX]
\baRequireAtlas
\end{lstlisting}

ensures that region declarations only occur inside a valid atlas environment.



\paragraph{Step 2: Region registry list}

\begin{lstlisting}[language=LaTeX]
\gappto\ba@regionlist{#2,}
\end{lstlisting}

This appends the region name to a global comma-separated list. This list is later used for iteration in rendering routines.

\paragraph{Step 3: Dynamic macro creation}

\begin{lstlisting}[language=LaTeX]
\expandafter\gdef\csname ba@caption@#2\endcsname{#1}
\end{lstlisting}

This constructs a macro dynamically:

\begin{itemize}
\item If region is \texttt{A}
\item Then macro becomes \texttt{\textbackslash ba@caption@A}
\end{itemize}

This mechanism allows storing region-specific metadata without predefined macros.



\paragraph{Step 4: TikZ style registration}

\begin{lstlisting}[language=LaTeX]
\tikzset{
  ba/region/#2/.style={
    /ba/internal/region colors={#3}
  }
}
\end{lstlisting}

Each region defines a TikZ style of the form:

\begin{lstlisting}[language=LaTeX]
ba/region/<name>
\end{lstlisting}

which internally stores the color configuration via the pgfkeys system.

\subsubsection{Appending Data with \texttt{\textbackslash gappto}}

The command

\begin{lstlisting}[language=LaTeX]
\gappto\ba@regionlist{#2,}
\end{lstlisting}

is used to build a growing, comma-separated list of region identifiers during compilation.

\paragraph{Meaning of \texttt{\textbackslash gappto}}

The macro \verb|\gappto| is provided by the \texttt{etoolbox} package. Its name stands for:

\begin{center}
\textbf{g}lobal \textbf{ap}pend \textbf{to}
\end{center}

It appends content to the end of an existing macro in a way that is \textbf{global}, meaning the change persists outside of local groups.



\paragraph{Operational behavior}

Given:

\begin{itemize}
\item existing macro: \verb|\ba@regionlist|
\item new content: \verb|#2,|
\end{itemize}

the command performs the equivalent of:

\begin{lstlisting}[language=LaTeX]
\xdef\ba@regionlist{\ba@regionlist#2,}
\end{lstlisting}

but in a safer and more controlled way.



\paragraph{Why the trailing comma is used}

The appended string:

\begin{lstlisting}[language=LaTeX]
#2,
\end{lstlisting}

ensures that the list becomes a comma-separated sequence:

\begin{lstlisting}[language=LaTeX]
A,B,C,
\end{lstlisting}

This structure is intentionally compatible with \verb|\foreach| loops in TikZ/PGF, which can directly iterate over comma-separated lists.



\paragraph{Why \texttt{gappto} is preferred over manual definitions}

Alternative approaches include:

\begin{lstlisting}[language=LaTeX]
\edef\ba@regionlist{\ba@regionlist#2,}
\end{lstlisting}

However, \verb|\gappto| is preferred because:

\begin{itemize}
\item it avoids low-level expansion pitfalls,
\item it preserves macro safety in complex environments,
\item it integrates well with other \texttt{etoolbox} mechanisms,
\item it makes intent (\emph{append globally}) explicit.
\end{itemize}



\paragraph{Role in the region system}

Within the bifurcation atlas package, \verb|\ba@regionlist| acts as a lightweight registry of all declared regions. Each call to:

\begin{lstlisting}[language=LaTeX]
\baDeclareRegion
\end{lstlisting}

extends this registry. The resulting list is later consumed by rendering routines such as:

\begin{itemize}
\item radial layout construction,
\item iteration over regions,
\item color assignment and caption mapping.
\end{itemize}

Thus, \verb|\gappto| provides the fundamental mechanism for building a dynamic, order-preserving data structure inside TeX.

% =========================================================
\subsubsection{Safe Access to Region Metadata}
% =========================================================

Region captions are retrieved using a conditional macro:

\begin{lstlisting}[language=LaTeX]
\newcommand{\baGetRegionCaption}[1]{%
  \ifcsname ba@caption@#1\endcsname
    \csname ba@caption@#1\endcsname
  \fi
}
\end{lstlisting}

\paragraph{Mechanism}

\begin{itemize}
\item \verb|\ifcsname ... \endcsname| checks whether a macro exists.
\item If it exists, it is expanded via \verb|\csname ... \endcsname|.
\item Otherwise, nothing is returned.
\end{itemize}

This avoids compilation errors when a region has no caption.



% =========================================================
\subsubsection{Color Resolution System}
% =========================================================

Colors are resolved via indexed macros:

\begin{lstlisting}[language=LaTeX]
\newcommand{\ba@getcolor}[1]{%
  \ifcsname ba@c#1\endcsname
    \csname ba@c#1\endcsname%
  \else white\fi
}
\end{lstlisting}

\paragraph{Idea}

Each color is stored as:

\begin{itemize}
\item \texttt{\textbackslash ba@c1}, \texttt{\textbackslash ba@c2}, ...
\end{itemize}

so that numeric indexing can be used instead of direct lists.



\paragraph{Public wrapper}

\begin{lstlisting}[language=LaTeX]
\newcommand{\baGetColor}[1]{\ba@getcolor{#1}}
\end{lstlisting}

This provides a stable public interface while keeping the internal implementation hidden.



% =========================================================
\subsubsection{Color List Expansion}
% =========================================================

Before rendering, the color list is expanded into indexed macros:

\begin{lstlisting}[language=LaTeX]
\newcommand{\baPrepareColors}{%
  \let\colors\ba@colors%
  \foreach [count=\i] \c in \colors {%
    \expandafter\xdef\csname ba@c\i\endcsname{\c}%
  }%
}
\end{lstlisting}

\paragraph{Step-by-step behavior}

\begin{itemize}
\item \verb|\let\colors\ba@colors| copies the stored list into a local iterable variable.
\item \verb|\foreach| iterates over each color in the list.
\item \verb|\xdef| globally defines indexed macros \verb|\ba@c1|, \verb|\ba@c2|, etc.
\end{itemize}

\paragraph{Result}

A list such as:

\begin{lstlisting}[language=LaTeX]
red,blue,green
\end{lstlisting}

becomes:

\begin{lstlisting}[language=LaTeX]
\ba@c1 = red
\ba@c2 = blue
\ba@c3 = green
\end{lstlisting}

This enables constant-time lookup during rendering.



\subsubsection{Design Summary}

This subsystem demonstrates a typical LaTeX package architecture:

\begin{itemize}
\item \textbf{Dynamic macros} for metadata storage
\item \textbf{pgfkeys integration} for structured configuration
\item \textbf{Indexed expansion} for efficient rendering
\item \textbf{Safe conditional access} to avoid undefined macro errors
\end{itemize}

Together, these mechanisms form a lightweight database-like layer inside TeX, enabling structured region management without external data structures.